\scanpage{6}

\begin{center}
U N I V E R S I T E\\
DES SCIENCES ET TECHNIQUES\\
DU LANGUEDOC\\
\textasciicircum\textasciicircum\textasciicircum\textasciicircum\textasciicircum\textasciicircum\textasciicircum\textasciicircum\textasciicircum\textasciicircum\textasciicircum\textasciicircum\textasciicircum
\end{center}

Montpellier, le 14 octobre 1983

Service de la Scolarité

Diffusion : Liste A\\
MA/gs \hfill N° 4375

\begin{center}
AVIS DE SOUTENANCE D'UNE THESE DE DOCTORAT DE 3e CYCLE

\uuline{ECOLOGIE GENERALE ET APPLIQUEE OPTION ECOLOGIE AQUATIQUE}
\end{center}

\uline{NOM et Prénom du candidat} : MONTI Dominique

\uline{Titre de la thèse} : ``Recherches sur l'optimisation du plan d'échantillonnage employé pour la cartographie des peuplements benthiques. Application à la macrofaune benthique de l'étang de Pierre Blanche (Palavas, Hérault, France).

\uline{Date,-heure et lieu de soutenance prévus} : Le 2 novembre 1983 à 15 heures,\\
Salle de cours n° 9 - Université des Sciences et Techniques du Languedoc

\uline{Nom du Directeur de thèse} : M. AMANIEU

\uline{Désignation du Laboratoire où la thèse est préparée} : Laboratoire d'Hydrobiologie marine - U.S.T.L.

\uline{Analyse de l'ouvrage} : Cette réflexion porte sur le problème du choix du nombre et de l'organisation spatiale des points de prélèvement en écologie. Quel effort doit-on fournir pour un niveau de précision fixé ? Et où doit-on prélever ?

Ce travail est basé sur la réalisation de différents sous-échantillonnages d'une information initiale supposée maximale. Les différents sous-échantillonnages varient par le nombre et l'organisation spatiale des points de prélèvements ; et ils sont tous effectués à partir de la même liste d'objets.

Chaque sous-échantillonnage employé donnera naissance à une carte qui sera par la suite comparée à la carte dite ``de référence'' ; il restera à apprécier au vu de ces résultats si les cartes obtenues diffèrent de la référence, en quelles proportions et de quelle manière.

Les résultats théoriques sont appliqués à l'étude de la macrofaune benthique d'un étang littoral (l'étang de Pierre Blanche à Palavas, France).

\uline{Durée approximative des recherches ayant abouti à la thèse} : 3 ans

\uline{Composition du jury Proposé} : \begin{tabular}[t]{ll}
MM. & AMANIEU M.\\
& BLANC F.\\
& FRONTIER S.\\
& GRANDJOUAN G.\\
& LASSERRE G.
\end{tabular}

\begin{flushright}
P/ le Secrétaire Général,\\
L'Attachée Principale d'Administration\\
Scolaire et Universitaire,

\unsure{A. Gentil}

A. GENTIL
\end{flushright}

\margin{5}
